\section{6. System Technical Budgets, Margins and
Deviations}\label{system-technical-budgets-margins-and-deviations}

\subsection{6.1 Budget Status and Sizing
Case}\label{budget-status-and-sizing-case}

\textbf{B} denotes a value from the existing requirements baseline, and
\textbf{P} a proposed design allocation that has not yet been approved
or demonstrated by testing. Calculations justifying the allocations and
design option comparisons belong in the GSEG DDJF.

\begingroup\GSEGTableStyle{3}
\begin{longtable}[]{|P{0.34}|P{0.39}|P{0.27}|}
\hline\rowcolor{TableHeaderBlue}
\begin{minipage}[b]{\linewidth}\raggedright
\color{white}\bfseries\sffamily
Parameter
\end{minipage} & \begin{minipage}[b]{\linewidth}\raggedleft
\color{white}\bfseries\sffamily
Value
\end{minipage} & \begin{minipage}[b]{\linewidth}\raggedright
\color{white}\bfseries\sffamily
Status and source
\end{minipage} \\
\hline
\endhead
\hline
\endlastfoot
Constellation & 6 GEO + 72 LEO & B; system/SSEG \\ \hline
Housekeeping telemetry generation rate & ≥1 sample/s per spacecraft & B;
\texttt{SSEG-\allowbreak{}REQ-\allowbreak{}OPS-\allowbreak{}006} \\ \hline
GSEG time intervals & ≤2 s / ≤15 s & B;
\texttt{GSEG-\allowbreak{}REQ-\allowbreak{}PERF-\allowbreak{}001/\allowbreak{}002} \\ \hline
Data retention & ≥5 years & B; \texttt{GSEG-\allowbreak{}REQ-\allowbreak{}DATA-\allowbreak{}005} \\ \hline
GSEG critical-service availability & ≥99.99\%; maximum
52.56 minutes unavailable per 365 days & B;
\texttt{GSEG-\allowbreak{}REQ-\allowbreak{}REL-\allowbreak{}003} \\ \hline
Concurrent active events & 10 & P; load scenario \\ \hline
Events in the archive model & 20 per day & P; volume model \\ \hline
GEO Event Report / LEO Tracking Request & 4 KiB / 4 KiB & P; maximum
modelled sizes \\ \hline
LEO Observation Package / preliminary trajectory & 2 MiB / 0.25 MiB & P;
maximum modelled sizes \\ \hline
Warning / addressees & 64 KiB / 10 & P; dissemination scenario \\ \hline
Archived housekeeping record & 2 KiB per sample & P; estimate \\ \hline
Indexing, audit and metadata overhead & 40\% of source volume & P;
planning factor \\ \hline
Hypothetical aggregate input from 20 LEO links & 3 Gb/s at 150 Mb/s per
link & P; load model, not demonstrated visibility \\*
\end{longtable}
\endgroup

The 1 Hz rate concerns \textbf{generation of telemetry on the
spacecraft}, not demonstrated continuous reception at GSN. The archive
model provisionally assumes that an equivalent volume is received; this
must be checked by SSEG/GSEG interface analysis. Ten simultaneous events
constitute a design scenario, not a binding upper limit on event
occurrence.

\newpage
\subsection{6.2 Timing Allocations}\label{timing-allocations}

\subsubsection{6.2.1 System Boundaries}\label{system-boundaries}

\begingroup\GSEGTableStyle{3}
\begin{longtable}[]{|P{0.62}|P{0.20}|P{0.18}|}
\hline\rowcolor{TableHeaderBlue}
\begin{minipage}[b]{\linewidth}\raggedright
\color{white}\bfseries\sffamily
Stage
\end{minipage} & \begin{minipage}[b]{\linewidth}\raggedright
\color{white}\bfseries\sffamily
Responsibility
\end{minipage} & \begin{minipage}[b]{\linewidth}\raggedleft
\color{white}\bfseries\sffamily
Limit
\end{minipage} \\
\hline
\endhead
\hline
\endlastfoot
GEO detection and reporting & SSEG & ≤10 s \\ \hline
LEO selection and request handover & \textbf{GSEG} & \textbf{≤2 s} \\ \hline
Command delivery to LEO & SSEG & ≤3 s \\ \hline
Observations and return of data to GSN & SSEG & ≤30 s \\ \hline
Ground assessment and warning delivery for warning-qualified events &
\textbf{GSEG} & \textbf{≤15 s} \\ \hline
\textbf{Entire EWAS: detection → delivery} & \textbf{System} &
\textbf{≤60 s} \\*
\end{longtable}
\endgroup

The allocations total 60 s, so the \textbf{unallocated system-level time
reserve is 0 s}. The GSEG reserves in the following tables are contained
within its limits. The measurement endpoint is delivery to the
authorized recipient interface, not a subsequent acknowledgement.

\subsubsection{6.2.2 Preliminary 2 s Budget
(P)}\label{preliminary-2-s-budget-p}

\begingroup\GSEGTableStyle{3}
\begin{longtable}[]{|P{0.63}|P{0.22}|P{0.15}|}
\hline\rowcolor{TableHeaderBlue}
\GSEGHeaderText{Operation} & \GSEGHeaderText{Owner} & \GSEGHeaderText{Allocation} \\
\hline
\endhead
\hline
\endlastfoot
Accept and make the report available & GSN & 0.20 s \\ \hline
GSN--C\&C transfer & GIS & 0.10 s \\ \hline
Select two LEO spacecraft and prepare requests & C\&C & 0.90 s \\ \hline
C\&C--GSN transfer & GIS & 0.10 s \\ \hline
Handover both requests to SSEG & GSN & 0.20 s \\ \hline
\textbf{Integration reserve} & GSEG & \textbf{0.50 s} \\ \hline
\textbf{Total} & & \textbf{2.00 s} \\*
\end{longtable}
\endgroup

C\&C uses orbital data previously prepared by MCC. GIS delays concern
internal exchange, not RF transmission to a spacecraft.

\subsubsection{6.2.3 Preliminary 15 s Budget
(P)}\label{preliminary-15-s-budget-p}

\begingroup\GSEGTableStyle{3}
\begin{longtable}[]{|P{0.63}|P{0.22}|P{0.15}|}
\hline\rowcolor{TableHeaderBlue}
\GSEGHeaderText{Operation} & \GSEGHeaderText{Owner} & \GSEGHeaderText{Allocation} \\
\hline
\endhead
\hline
\endlastfoot
Receive and hand over complete products & GSN & 0.40 s \\ \hline
Inter-service transfer & GIS & 0.40 s \\ \hline
Validate observation products & DPS & 1.00 s \\ \hline
Classification, fusion and impact-area estimation & DPS & 7.20 s \\ \hline
Prepare and version the warning & EWS & 1.00 s \\ \hline
Deliver to the authorized recipient interface & EWS & 2.00 s \\ \hline
Critical scheduling and queuing & GIS & 0.50 s \\ \hline
\textbf{Integration reserve} & GSEG & \textbf{2.50 s} \\ \hline
\textbf{Total} & & \textbf{15.00 s} \\*
\end{longtable}
\endgroup

The limits cover the ten-event scenario and failures covered by the
adopted redundancy (P). Failover does not add time beyond 15 s.

Applicability: the 15 s allocation covers assessment and
first-warning delivery \textbf{for events qualified by DPS for warning}.
A credibly non-threatening outcome or ``assessment impossible'' status
results in a recorded decision and an MCC notification, not an attempted
warning transmission; this is not a warning-delivery timing case.

\subsection{6.3 Technical Resource Allocations}\label{technical-resource-allocations-corrected-version}

GIS resource sizing (P): C\&C and EWS have
fixed per-domain resource envelopes shared across multiple events; DPS
demand scales with the number of active events. The capacity of
\emph{each} GIS domain covers the scenario in which the other domain has
failed.

\begingroup\GSEGTableStyle{3}
\begin{longtable}[]{|P{0.27}|P{0.39}|P{0.34}|}
\hline\rowcolor{TableHeaderBlue}
\begin{minipage}[b]{\linewidth}\raggedright
\color{white}\bfseries\sffamily
Resource / owner
\end{minipage} & \begin{minipage}[b]{\linewidth}\raggedright
\color{white}\bfseries\sffamily
Proposed design allocation P
\end{minipage} & \begin{minipage}[b]{\linewidth}\raggedright
\color{white}\bfseries\sffamily
Boundary / note
\end{minipage} \\
\hline
\endhead
\hline
\endlastfoot
GIS critical network & ≥10 Gb/s in each active domain & Does not
demonstrate GSN--SSEG RF link availability. \\ \hline
EWS warning egress & ≥100 Mb/s effective aggregate capacity for 10
recipients & Proposal; protocol overhead and availability unverified. \\ \hline
GIS critical platform & 2 separate domains; ≥160 vCPU
and ≥320 GiB RAM in each domain & Each domain assumes the entire
critical load after loss of the other. \\ \hline
C\&C --- shared service & 8 vCPU / 16 GiB RAM per
domain & \emph{Fixed} envelope for the ten-event scenario. \\ \hline
EWS --- shared service & 8 vCPU / 16 GiB RAM per
domain & \emph{Fixed} envelope for the ten-event scenario. \\ \hline
DPS --- event processing & 8 vCPU / 16 GiB RAM per
event × 10 = 80 vCPU / 160 GiB per domain & Linear load scaling in the
model; algorithm performance not demonstrated. \\ \hline
\textbf{Total C\&C + DPS + EWS} & \textbf{96 vCPU / 192 GiB per domain}
& Sum of critical-service allocations. \\ \hline
GIS platform and system-service overhead & 24 vCPU / 48 GiB per
domain & Planning hypothesis for runtime, monitoring and platform
overhead. \\ \hline
\textbf{Modelled nominal capacity in use} & \textbf{120 vCPU / 240 GiB
per domain} & 96+24 vCPU; 192+48 GiB. \\ \hline
Reserve after allocation & 40 vCPU / 80 GiB per domain
& 25\% of domain capacity; an allocation, not a demonstrated performance
reserve. \\ \hline
DAS archive & ≥50 TiB usable per copy; 2 independently maintained
logical copies & Both copies at one site; hardware overhead excluded
from usable capacity. \\ \hline
GIS critical IT power & 12 kW planning load & Excludes GSN RF, DAS, MCC
and whole-site utilities. \\ \hline
GIS critical IT cooling & ≥15 kW heat-removal capacity & Only for the
above 12 kW. \\ \hline
GIS critical IT UPS & ≥5 kWh usable energy for a 15-minute bridging case
& Not a whole-site power supply. \\*
\end{longtable}
\endgroup

These service envelopes do not prescribe whether shared services run as
single instances, clusters or multiple processes. The details belong to
subsystem DDFs. The 40 vCPU / 80 GiB reserve \textbf{does not replace a
timeliness test} for the 2 s and 15 s intervals.

\subsubsection{6.3.1 Preliminary GSEG Critical-Service Availability
Budget}\label{preliminary-gseg-critical-service-availability-budget}

The operational availability requirement for critical GSEG services is
≥99.99\% (\texttt{GSEG-\allowbreak{}REQ-\allowbreak{}REL-\allowbreak{}003}, status B). Over a continuous
365-day window this means at most 100 ppm unavailability, or
52.56 minutes in total. This is the requirement limit, not a
measured result.

Readiness definition: GSEG is available when \emph{both}
segment-controlled capabilities can operate within their deadlines upon
receipt of valid input data:

\begin{itemize}
\tightlist
\item
  \textbf{Tasking:} reception of a GEO report at GSN → handover of two
  requests to SSEG within \textbf{≤2 s}.
\item
  \textbf{Assessment and warning:} reception of valid products from two
  LEO spacecraft and the onboard preliminary trajectory at GSN →
  preliminary warning delivery within \textbf{≤15 s}, if DPS qualifies
  the event for warning.
\end{itemize}

The measurement covers continuous 24/7 readiness in a rolling 365-day
window. Unavailability includes planned maintenance, faults and
failover delays preventing either capability. In the absence of actual
events, readiness may be checked with controlled test stimuli.
Independent SSEG failures and failures of recipient systems beyond the
interface boundary are excluded; links and elements controlled by GSEG
are included.

\begingroup\GSEGTableStyle{3}
\begin{longtable}[]{|P{0.40}|P{0.20}|P{0.40}|}
\hline\rowcolor{TableHeaderBlue}
\begin{minipage}[b]{\linewidth}\raggedright
\color{white}\bfseries\sffamily
Unavailability component
\end{minipage} & \begin{minipage}[b]{\linewidth}\raggedleft
\color{white}\bfseries\sffamily
Budget P
\end{minipage} & \begin{minipage}[b]{\linewidth}\raggedleft
\color{white}\bfseries\sffamily
Target design availability
\end{minipage} \\
\hline
\endhead
\hline
\endlastfoot
\textbf{GSN} --- ground satellite communications & 20 ppm &
≥99.9980\% \\ \hline
\textbf{C\&C} --- request and command handling & 10 ppm & ≥99.9990\% \\ \hline
\textbf{DPS} --- event assessment & 10 ppm & ≥99.9990\% \\ \hline
\textbf{EWS} --- warning delivery & 15 ppm & ≥99.9985\% \\ \hline
\textbf{GIS} --- shared infrastructure and site dependencies & 25 ppm &
≥99.9975\% \\ \hline
\textbf{Integration and egress connections controlled by GSEG} & 10 ppm
& ≥99.9990\% \\ \hline
\textbf{Total allocated} & \textbf{90 ppm} & Allocation, not a measured
result \\ \hline
\textbf{Unallocated reserve} & \textbf{10 ppm} & Reserve within the GSEG
limit \\ \hline
\textbf{Total limit} & \textbf{100 ppm} & \textbf{≥99.9900\%} \\*
\end{longtable}
\endgroup

\begin{figure}[H]
\centering
% Joint service readiness and shared dependencies; no packet-flow arrows.
\resizebox{0.98\linewidth}{!}{%
\begin{tikzpicture}[
  service/.style={draw=HeadingBlue,fill=TableRowLightBlue,rounded corners=2pt,
    align=center,font=\small,text width=3.0cm,minimum height=1.25cm,inner sep=3pt},
  shared/.style={draw=TableBorderBlue,fill=white,rounded corners=2pt,
    align=center,font=\small,text width=12.8cm,minimum height=1.25cm,inner sep=4pt},
  readiness/.style={draw=HeadingBlue,line width=0.75pt},
  dependency/.style={draw=HeadingBlue,dashed,line width=0.65pt}
]
\node[service] (gsn) at (-5.7,0) {GSN\\Ground RF / gateway};
\node[service] (cc) at (-1.9,0) {C\&C\\Tasking service};
\node[service] (dps) at (1.9,0) {DPS\\Assessment service};
\node[service] (ews) at (5.7,0) {EWS\\Recipient handover};
\draw[readiness] (gsn.east) -- (cc.west);
\draw[readiness] (cc.east) -- (dps.west);
\draw[readiness] (dps.east) -- (ews.west);
\node[shared] (gis) at (0,-2.6)
  {GIS: shared compute, networks, time and site utilities\\
   Two on-site equipment domains; site-wide common-cause exposure remains};
\draw[dependency] (gis.north) -- (0,-1.4);
\draw[dependency] (-5.7,-1.4) -- (5.7,-1.4);
\foreach \name/\x in {gsn/-5.7,cc/-1.9,dps/1.9,ews/5.7} {
  \draw[dependency] (\x,-1.4) -- (\name.south);
}
\node[align=center,font=\small,text width=15.0cm] at (0,-4.9)
  {Both GSEG-controlled capabilities must be ready: $A_{\mathrm{GSEG}}\geq99.99\%$.\\
   Controlled integration and egress are included. Equipment-domain redundancy does not establish site independence.\\
   Lines show logical availability dependencies, not data flow or independent failures.};
\end{tikzpicture}%
}

\caption{Figure 6-1. Logical GSEG availability block model: joint
readiness of the GSN, C\&C, DPS and EWS services and their dependence on
GIS and segment-controlled connections. This is a service RBD, not a
physical-wiring diagram or proof of failure independence.}
\end{figure}

The preliminary Reliability Block Diagram (RBD) represents
jointly required critical-service capabilities. Allocations are
conservatively apportioned among unavailability components and reserve;
P values are not calculated MTBF/MTTR results or demonstrated
performance. The DDJF analysis must cover site common-cause failures,
GIS inter-domain dependencies, loss of GSN and outbound links, fault
detection, failover, repair and scheduled maintenance. Two GIS domains
alone do not demonstrate compliance with 99.99\%.

\newpage
\subsection{6.4 Margins and Deviations}\label{margins-and-deviations-corrected-version}

\begingroup\GSEGTableStyle{3}
\begin{longtable}[]{|P{0.30}|P{0.36}|P{0.34}|}
\hline\rowcolor{TableHeaderBlue}
\begin{minipage}[b]{\linewidth}\raggedright
\color{white}\bfseries\sffamily
Parameter
\end{minipage} & \begin{minipage}[b]{\linewidth}\raggedright
\color{white}\bfseries\sffamily
Design margin / reserve
\end{minipage} & \begin{minipage}[b]{\linewidth}\raggedright
\color{white}\bfseries\sffamily
Evidence status
\end{minipage} \\
\hline
\endhead
\hline
\endlastfoot
GSEG 2 s interval & 0.50 s = 25\% of limit & P; allocation reserve, not
measured. \\ \hline
GSEG 15 s interval & 2.50 s = 16.7\% of limit & P; allocation reserve,
not measured. \\ \hline
Entire EWAS 60 s & 0 s unallocated & B; sum of system allocations. \\ \hline
GSEG critical-service availability & ≥99.99\%;
100 ppm limit with 90 ppm allocated and \textbf{10
ppm} reserve & Adopted requirement; P allocation without verified
analysis or measurement. \\ \hline
GIS critical network & Theoretically 70\% unused capacity at 3 Gb/s out
of 10 Gb/s & P; visibility, overhead and contact windows unverified. \\ \hline
Compute resources in surviving GIS domain & 40 vCPU /
80 GiB = 25\% of 160 vCPU / 320 GiB & P; after subtracting 96 vCPU /
192 GiB services and 24 vCPU / 48 GiB overhead; not measured performance
headroom. \\ \hline
DAS five-year archive & Approximately 35\% of each 50 TiB copy remains
free & P; depends on the data-stream and event-count model. \\ \hline
Critical IT cooling & 20\% of the 15 kW capacity unused under a 12 kW
load & P; not a whole-site margin. \\ \hline
Verified GSEG margins & No available results & No confirmed integrated
analysis or test has been completed. \\ \hline
Confirmed deviations from GSEG TRS & None formally identified & Does not
imply demonstrated compliance. \\*
\end{longtable}
\endgroup

The P values are \textbf{proposed allocations}, not approved
requirements. Only subsequent load, satellite-contact, latency and
failover analysis can establish resource adequacy. Selection methods and
option comparisons belong in the DDJF. A confirmed negative margin
requires controlled assessment and corrective action.
